\hypertarget{a-java-discovery-boundary-that-retains-the-question}{%
\section{A Java discovery boundary that retains the
question}\label{a-java-discovery-boundary-that-retains-the-question}}

\label{exec13:java-adapter}

The Java comparison motivates a concrete first integration: a pinned
Tetrad process for a restricted PC discovery request, with a separate
admission path for FGES. The initial request concerns complete
continuous observations under an explicitly assumed IID linear Gaussian
model with causal sufficiency, acyclicity, faithfulness and no unmodeled
selection. These restrictions identify a reviewable prototype. They do
not imply that a table of banking observations satisfies them. Clustered
customers, repeated account histories, missing-data correction,
latent-variable discovery and background-knowledge constraints require
other profiles. The adapter is a proposal contract; no Java engine or
new transport endpoint has been implemented here.
\citep{exec13_java_adapter_contract,exec13_java_tetrad}

The request retains the downstream target even though discovery does not
compute it. A funding analyst's population, treatment contrast, outcome,
time horizon, units and sampling declarations remain attached to the
graph. They prevent a later identification step from silently changing
the question. The generic operation contract already provides these
fields for effect operations. The Java discovery profile deliberately
requires that context as well. A standalone graph request without a
downstream treatment and outcome needs its own reviewed context
extension; inserting arbitrary variable names would create a fictitious
causal question. \citep{exec13_java_adapter_contract}

The logical interface is the same across the proposed library
invocation, \texttt{causalbank operation graph.discover} command and
\texttt{causal\_graph\_discover} MCP tool. Each transports a versioned
canonical request and receives the same typed result. A complete
scientific context digest accompanies a broader request digest
containing input hashes, engine revision, algorithm settings, adapter
version and resource limits. An interface that changes USD into
thousands of USD, reverses the treatment contrast, changes the sampling
law or orients an ambiguous edge has changed the scientific request,
even if its JSON remains valid. The reference serialization examples
demonstrate those identity changes on local objects; they are not
executions of a server or MCP client.
\citep{exec13_java_adapter_contract,exec13_java_adapter_oracles}

\hypertarget{the-graph-object-preserves-a-class-of-possible-interventions}{%
\subsection{The graph object preserves a class of possible
interventions}\label{the-graph-object-preserves-a-class-of-possible-interventions}}

Use opaque variable identifiers and a bijection into engine symbols.
Store the original column order and its permutation separately from the
engine's order. Every variable, including an isolated node, travels
through the conversion. Human-readable labels, availability times and
units remain metadata with their own scientific meaning; names do not
supply an orientation rule. For every pair, retain its two endpoint
marks. When the serialization orders the pair lexically, it swaps the
marks with the nodes. Thus \(R\to B\), expressed as a tail at \(R\) and
an arrowhead at \(B\), agrees with the reversed pair representation
having an arrowhead at \(B\) and a tail at \(R\). Reversing only the
node order changes the graph.
\citep{exec13_java_tetrad,exec13_java_adapter_oracles}

A Gaussian example illustrates why an unoriented link must survive this
boundary. Consider three standardized variables \(X,Y,Z\) with
covariance \[
 \Sigma=\begin{pmatrix}
 1&4/5&2/5\\
 4/5&1&1/2\\
 2/5&1/2&1
 \end{pmatrix}.
\] For a banking interpretation, \(X\) can denote a rate concession,
\(Y\) a balance response and \(Z\) a funding outcome, all in stipulated
standardized units. The covariance obeys \(X\perp Z\mid Y\). The
undirected path \(X-Y-Z\) represents three acyclic orientations without
an unshielded collider. The orientation \(X\to Y\leftarrow Z\) is
excluded because it would require a different independence pattern under
faithfulness. \citep{exec13_java_adapter_oracles}

For each retained orientation, regress a variable on its proposed
parents using the same \(\Sigma\), and assign the residual variance as
its independent Gaussian noise variance. In matrix notation, if
\(W_{ij}\) is the coefficient from source \(i\) to target \(j\), then \[
 \Sigma=(I-W^{\mathsf T})^{-1}\Omega(I-W)^{-1},
\] where \(\Omega\) is the diagonal noise covariance. The exact
reference calculation verifies this equality for all three models. In
the forward chain, the coefficients are \(4/5\) and \(1/2\); intervening
on \(X\) changes \(Z\) by \(2/5\) per unit. In the fork and the reverse
chain, \(X\) has no directed path to \(Z\), so the effect is zero. The
same observational covariance therefore supports the effect multiset
\(\{2/5,0,0\}\). \citep{exec13_java_adapter_oracles}

Returning the forward chain alone would select the positive effect by an
unrecorded convention. Returning the interval \([0,2/5]\) would
introduce values such as \(1/5\) that none of these three specified
linear Gaussian models produces. That interval is a hull, not the
discrete effect set, and neither object is a confidence set without a
separate statistical construction. The discovery response consequently
retains the CPDAG and returns no effect estimate, effect region or
uncertainty level. A later identification operation must state how it
handles the whole class and the additional functional assumptions.
\citep{exec13_java_adapter_contract,exec13_java_adapter_oracles}

Checking endpoint syntax is insufficient to award the CPDAG label. A
graph can contain only tails and arrows and still fail to represent a
completed Markov equivalence class. The existing reference validator
does not recognize all such cases. The proposed Java boundary therefore
requires an independently reviewed graph-class check before publishing a
completed CPDAG. A circle, bidirected conflict or invalid completion in
the initial route is an invalid-engine-output failure, with the raw
graph retained for diagnosis. It is not silently coerced into a DAG.
Background knowledge is initially excluded because orienting an
observational CPDAG using additional restrictions requires an
appropriate constrained-graph representation and provenance.
\citep{exec13_java_adapter_contract}

\hypertarget{admission-must-cover-the-failure-path-as-well-as-the-constructor}{%
\subsection{Admission must cover the failure path as well as the
constructor}\label{admission-must-cover-the-failure-path-as-well-as-the-constructor}}

The proposed PC profile sets the conditioning depth to the inspected
unbounded value, enables the stable skeleton option, selects the
separating-set collider rule, disallows bidirected conflicts, and
prevents directed cycles. It supplies empty background knowledge and a
caller-declared test level. These choices make the attempted operation
explicit. A user-requested depth cap is initially unsupported;
exhausting a runtime budget returns an incomplete job, not a graph
advertised as the result of the completed search.
\citep{exec13_java_adapter_contract,exec13_java_adapter_source}

The numerical test profile is equally specific: complete data, no
shrinkage, zero ridge and no pseudoinverse fallback, with all settings
frozen before search. The inspected Fisher-Z source exposes these
alternatives separately. Automatically activating regularization after a
singularity would change the test while keeping the same apparent method
name. That may be useful in a different profile, but it needs a separate
calibration argument. The initial profile requires at least
\(n\geq p+2\) complete observations for \(p\) variables so that even the
largest possible conditioning set has positive Fisher-Z degrees of
freedom. This is only an arithmetic prerequisite; it does not guarantee
useful power or stable covariance inversion.
\citep{exec13_java_adapter_contract,exec13_java_adapter_source}

As the Java comparison showed, the inspected \texttt{checkIndependence}
path converts a singular-matrix exception into a dependence result with
\(p=0\). The proposed wrapper instead uses the exception-preserving
\texttt{getResult} path and derives a decision only from a valid result.
Every invocation receives a query identity and a terminal record. Before
propagating an exception, the wrapper sets a job-level failure flag.
Publication then checks that every issued call has a valid terminal
record and that the flag is clear. This matters because search code may
catch interruptions internally: receiving a graph object alone does not
prove that all its requested tests completed.
\citep{exec13_java_tetrad,exec13_java_adapter_source,exec13_java_adapter_contract}

That wrapper remains an implementation prerequisite. Its full
independence-test interface, ancillary methods and interaction with the
search implementation must be inspected and compiled; calling one method
in a design document is not sufficient. Numerical rank checks before
dispatch help detect duplicate accounting columns, but per-query records
remain necessary for failures that arise during the search. A zero test
probability cannot be reinterpreted as maximal causal confidence. If the
required failure provenance cannot be obtained, the route is unavailable
in this runtime. \citep{exec13_java_adapter_contract}

Parallelism and deadlines also belong to admission. The stable FAS path
uses a parallel stream and requires thread-safe records. FGES's separate
worker setting does not control that path. The inspected PC search
delegates the skeleton phase to FAS without demonstrating an end-to-end
deadline guarantee from PC's timeout setting alone. The proposed parent
process therefore enforces the whole-job deadline, memory and output
budget. An engine's cooperative timeout is supplementary. A cancellation
or a process death preserves the attempt's failure record and withholds
its graph.
\citep{exec13_java_adapter_source,exec13_java_adapter_contract}

\hypertarget{units-unsupported-requests-and-the-next-executable-comparison}{%
\subsection{Units, unsupported requests and the next executable
comparison}\label{units-unsupported-requests-and-the-next-executable-comparison}}

The signed funding model in the Java survey provides a useful
independent conversion reference. If \(x'_i=d_i x_i\), the source-row
coefficients must transform as \(W'_{ij}=d_jW_{ij}/d_i\). Rescaling the
rate by 100, balances by 1000 and cost by 10, then applying the
correspondingly rescaled rate intervention, must recover the original
responses after converting units back. The exact reference does so. It
also verifies that an unchanged absolute coefficient threshold can
select a different graph after rescaling. This does not give
coefficients to the PC adapter, whose result is graph-only; it supplies
an admission example for any later weighted route.
\citep{exec13_java_adapter_oracles,exec13_java_examples}

A valid request outside the initial scientific scope receives a
completed unsupported response with its context preserved and a specific
alternative. Repeated account observations point toward a separately
reviewed temporal or cluster profile; latent-variable assumptions point
toward an appropriate PAG route; missing values require an explicit
observation mechanism and correction profile. An uninstalled or
unqualified runtime receives an unavailable response. Neither condition
is called nonidentification, which would be a stronger scientific
conclusion about an effect under a model. Malformed inputs and failed
computations carry structured errors, with no effect or graph result
masquerading as a completed calculation.
\citep{exec13_java_adapter_contract}

The first prototype should use a fixed local Java process so its
dependency and runtime identity can be measured. The process boundary
does not remove the inspected GPL conditions. The actual compiler
target, runtime vendor and version, dependency artifacts, operating
system and thread configuration must be locked for the chosen route; the
source's compiler setting and bridge README cannot substitute for that
environment check. Result publication is atomic after validation. A
crashed search is restarted as a new retained attempt under the same
immutable request, with no promise of resuming an internal search state.
If completion wins a cancellation race, the immutable published result
remains the observed outcome.
\citep{exec13_java_tetrad,exec13_java_adapter_contract}

The executable comparison begins with independent oracles: the
three-model Gaussian class, a collider and an isolated node, the
depth-three separator, signed matrices and unit changes, singular input,
and injected failures. Only after these pass should a prospectively
specified generated-sample study estimate graph error and resource use
under the declared assumptions and their violations. It should retain
every attempted request, including timeouts and refusals, and compare a
matched baseline under the same budgets. These steps test whether the
proposed boundary deserves a prototype claim. Expensive bank experiments
become relevant when that bounded evidence exposes a decision-relevant
uncertainty that public or generated data cannot settle.
\label{exec13:java-adapter-end}

\hypertarget{encoding-algorithm-settings-without-changing-the-shared-schema}{%
\subsection{Encoding algorithm settings without changing the shared
schema}\label{encoding-algorithm-settings-without-changing-the-shared-schema}}

The existing discovery request schema permits only an algorithm
identifier and algorithm version in its parameter object. The proposed
Java profile therefore places the detailed settings in a separate
immutable configuration artifact. The request supplies its opaque
handle, hash, media type, schema version and role, alongside the data
artifact. Resolving that artifact recovers the test level, conditioning
depth, missingness policy, numerical controls and background-knowledge
choice. The full request identity includes its hash, so changing a
setting creates a different request even when the downstream scientific
context remains identical. This preserves compatibility with the shared
schema while retaining the scientifically relevant settings.
\citep{exec13_java_adapter_contract}

The adapter's admission check must require exactly one configuration
artifact of the reviewed version, verify its bytes and recognized
fields, and check that the declared algorithm version admits that
profile. Generic schema validation alone cannot establish those
conditions. The worked wire records separately represent an unavailable
runtime, an unsupported sampling regime, a failed invocation and a graph
response from an independent population reference. Their data fixture is
deterministic, and their graph is not a fitted result. These records
make the proposed response contract concrete without treating a valid
message as evidence of an executed engine. \label{exec13:java-wire-end}
